% Chapitre — État de l'art (insertion après ch.3)
\chapter{État de l'art et fondements théoriques}
\label{ch:etatart}

\section{Scoring crédit : approches traditionnelles}

Le \textbf{scoring crédit} consiste à estimer la probabilité qu'un emprunteur ne rembourse pas ses obligations. Historiquement, deux grandes familles coexistent :

\subsection{Modèles statistiques paramétriques}

La \textbf{régression logistique} modélise :
\[
P(Y=1 \mid X) = \frac{1}{1 + e^{-(\beta_0 + \beta^T X)}}
\]
où $Y$ est la variable défaut et $X$ le vecteur de covariables. Avantages : interprétabilité des coefficients, rapidité, baseline robuste. Limites : hypothèse de linéarité du log-odds.

\subsection{Modèles d'ensemble}

Les \textbf{forêts aléatoires} (Breiman, 2001) agrègent des arbres de décision entraînés sur des sous-échantillons bootstrap. Elles capturent les non-linéarités et interactions sans feature engineering explicite.

\textbf{XGBoost} (Chen \& Guestrin, 2016) optimise un gradient boosting régularisé, dominant de nombreux concours Kaggle tabulaires. Hyperparamètres clés : profondeur d'arbres, learning rate, nombre d'estimateurs.

\section{Scoring en microfinance}

La microfinance présente des spécificités :
\begin{itemize}[leftmargin=*]
  \item montants faibles, fréquence élevée de remboursements ;
  \item absence d'historique bancaire formel (thin-file borrowers) ;
  \item importance du réseau social (group lending, garants) ;
  \item contraintes réglementaires KYC/AML renforcées.
\end{itemize}

Les institutions s'appuient sur des \textbf{scorecards} combinant revenu, activité, garanties et comportement de remboursement. L'IA permet d'intégrer des signaux alternatifs (transactions mobile money, géolocalisation — hors périmètre de ce PFE).

\section{Modèles séquentiels pour le comportement}

Les réseaux de neurones récurrents (RNN) modélisent des \textbf{séries temporelles} de transactions :
\begin{itemize}[leftmargin=*]
  \item \textbf{LSTM} (Hochreiter \& Schmidhuber, 1997) : portes forget/input/output, capture dépendances longues ;
  \item \textbf{GRU} (Cho et al., 2014) : architecture simplifiée, moins de paramètres, entraînement plus rapide.
\end{itemize}

Application crédit : détecter une dégradation progressive (retards croissants, baisse des dépôts) avant le défaut effectif.

\section{Graph Neural Networks}

Les \textbf{GNN} généralisent le deep learning aux données graphe $G=(V,E)$. \textbf{GraphSAGE} (Hamilton et al., 2017) apprend des représentations par échantillonnage de voisinages :
\[
\mathbf{h}_v^{(k)} = \sigma\left(W^{(k)} \cdot \text{AGG}\left(\mathbf{h}_v^{(k-1)}, \{\mathbf{h}_u^{(k-1)} : u \in \mathcal{N}(v)\}\right)\right)
\]

En microfinance, le graphe clients-relations modélise la \textbf{contagion de risque} : un défaut de garant peut signaler un risque sur le client principal.

\section{Explainable AI (XAI)}

La réglementation (GDPR art. 22, directives Bâle) exige une \textbf{explicabilité} des décisions automatisées. Approches :
\begin{itemize}[leftmargin=*]
  \item \textbf{Intrinsèque} : régression logistique, arbres ;
  \item \textbf{Post-hoc} : SHAP, LIME sur modèles boîte noire ;
  \item \textbf{LLM} : génération de texte à partir de features structurées (notre approche).
\end{itemize}

Notre projet combine scores ML déterministes + explication LLM + citations RAG pour traçabilité documentaire.

\section{Retrieval-Augmented Generation}

Lewis et al. (2020) formalisent le \textbf{RAG} : un retriever $p_\eta(z|x)$ sélectionne des documents $z$, un generator $p_\theta(y|x,z)$ produit la réponse $y$. Avantages :
\begin{itemize}[leftmargin=*]
  \item réduction des hallucinations ;
  \item mise à jour du knowledge sans retraining LLM ;
  \item citations vérifiables.
\end{itemize}

Notre implémentation TF-IDF est un baseline ; les embeddings denses (BERT, E5) améliorent la sémantique.

\section{Agents LLM et LangGraph}

Les \textbf{agents} LLM utilisent des \textit{tools} (appels API, bases de données) dans une boucle planification-action. LangGraph structure cette boucle en \textbf{graphe d'états fini}, offrant :
\begin{itemize}[leftmargin=*]
  \item reproductibilité des workflows ;
  \item checkpoints et reprise ;
  \item parallélisation de branches ;
  \item séparation claire des responsabilités (nœuds SRP).
\end{itemize}

\section{Positionnement de notre contribution}

\begin{table}[H]
  \centering
  \caption{Contribution vs littérature}
  \begin{tabular}{|p{4cm}|p{5cm}|p{4cm}|}
    \hline
    \textbf{Aspect} & \textbf{Littérature} & \textbf{Notre projet} \\
    \hline
    Multi-paradigme & Souvent un seul modèle & 3 familles orchestrées \\
    \hline
    Interaction & API ou batch & Chat LangGraph \\
    \hline
    Explicabilité & SHAP ou texte libre & LLM + RAG REF-08 \\
    \hline
    Déploiement & Cloud ML & Local Windows + Ollama \\
    \hline
  \end{tabular}
\end{table}

\section{Normes et conformité}

\begin{itemize}[leftmargin=*]
  \item \textbf{Bâle II/III} : exigences fonds propres selon risque crédit ;
  \item \textbf{KYC/AML} : identification client, lutte blanchiment ;
  \item \textbf{Protection des données} : minimisation, finalité, droit d'accès (RGPD-like).
\end{itemize}

Le prototype inclut des disclaimers et sépare scoring automatique de décision humaine, préparant une certification future.
